\documentclass[11pt,a4paper]{article}

\usepackage[utf8]{inputenc}
\usepackage[T1]{fontenc}
\usepackage{geometry}
\geometry{margin=2.5cm}
\usepackage{listings}
\usepackage{xcolor}
\usepackage{hyperref}
\usepackage{enumitem}
\usepackage{booktabs}
\usepackage{parskip}
\usepackage{titlesec}

\hypersetup{colorlinks=true, linkcolor=blue!50!black, urlcolor=blue!50!black, citecolor=blue!50!black}

\definecolor{codebg}{RGB}{245,245,245}
\definecolor{codegray}{RGB}{110,110,110}
\definecolor{codegreen}{RGB}{0,110,0}
\definecolor{codeblue}{RGB}{30,80,180}

\lstdefinestyle{code}{
    backgroundcolor=\color{codebg}, basicstyle=\ttfamily\footnotesize,
    keywordstyle=\color{codeblue}, commentstyle=\color{codegray}\itshape,
    stringstyle=\color{codegreen}, breaklines=true, breakatwhitespace=false,
    frame=single, rulecolor=\color{black!15}, numbers=left,
    numberstyle=\tiny\color{codegray}, xleftmargin=1.8em, framexleftmargin=1.5em,
    showstringspaces=false, tabsize=2,
    literate=
      {á}{{\'a}}1 {à}{{\`a}}1 {ã}{{\~a}}1 {â}{{\^a}}1
      {é}{{\'e}}1 {ê}{{\^e}}1 {í}{{\'i}}1
      {ó}{{\'o}}1 {õ}{{\~o}}1 {ô}{{\^o}}1 {ú}{{\'u}}1
      {ç}{{\c c}}1 {Á}{{\'A}}1 {É}{{\'E}}1 {Í}{{\'I}}1
      {Ó}{{\'O}}1 {Ú}{{\'U}}1 {Ã}{{\~A}}1 {Õ}{{\~O}}1
      {Ç}{{\c C}}1,
}
\lstset{style=code}

\titleformat{\section}{\normalfont\Large\bfseries}{\thesection}{1em}{}
\titleformat{\subsection}{\normalfont\large\bfseries}{\thesubsection}{1em}{}

\sloppy

\title{\textbf{TelemetriaMqtt v2 --- Checkpoint 5} \\[4pt]
\large Schema e setup do InfluxDB}
\author{Projeto de Telemetria FSAE}
\date{16 de setembro de 2026}

\begin{document}
\maketitle

\begin{abstract}
\noindent Este documento detalha o quinto checkpoint do projeto: colocar o InfluxDB no ar como container Docker, confirmar de fato a suposição de versão (2.x) feita no checkpoint 4, e validar contra um banco real -- não mais mockado -- a escrita e a leitura de pontos de telemetria.
\end{abstract}

\tableofcontents

\section{Contexto e objetivo}

No checkpoint 4, o script de ingestão foi escrito e testado com um \texttt{write\_api} \textbf{mockado} (\texttt{unittest.mock.MagicMock}), porque o InfluxDB ainda não existia como serviço. A suposição de que o banco seria InfluxDB 2.x (baseada na terminologia de "bucket" usada nos slides do TCC antigo) ficou registrada no \texttt{SCOPE.md} como pendente de confirmação. Este checkpoint resolve essa pendência: coloca um InfluxDB 2.x real no ar, e valida contra ele.

\section{Estrutura adicionada}

\begin{lstlisting}[language=bash,caption={Arquivos novos/alterados}]
infra/
  docker-compose.yml        (+ serviço influxdb)
  .env.example               (+ variáveis do InfluxDB)
  requirements-test.txt      (+ influxdb-client)
  tests/
    conftest.py               (+ fixture ensure_env_file, influxdb_instance)
    test_influxdb.py           (novo)
    test_docker_compose.py     (+ teste do serviço influxdb)
    test_end_to_end.py         (+ teste do pipeline completo)
\end{lstlisting}

\subsection{Serviço no docker-compose.yml}

\begin{lstlisting}[language=yaml,caption={Trecho de infra/docker-compose.yml}]
influxdb:
  image: influxdb:2
  environment:
    DOCKER_INFLUXDB_INIT_MODE: setup
    DOCKER_INFLUXDB_INIT_USERNAME: ${INFLUXDB_INIT_USERNAME:-admin}
    DOCKER_INFLUXDB_INIT_PASSWORD: ${INFLUXDB_INIT_PASSWORD:?...}
    DOCKER_INFLUXDB_INIT_ORG: ${INFLUXDB_INIT_ORG:-fsae}
    DOCKER_INFLUXDB_INIT_BUCKET: ${INFLUXDB_INIT_BUCKET:-telemetria}
    DOCKER_INFLUXDB_INIT_ADMIN_TOKEN: ${INFLUXDB_INIT_ADMIN_TOKEN:?...}
  volumes:
    - influxdb-data:/var/lib/influxdb2
\end{lstlisting}

A imagem oficial faz o setup inicial (usuário, organização, bucket, token de admin) sozinha, lendo variáveis de ambiente, na primeira subida com o volume de dados vazio.

\subsection{Credenciais via infra/.env}

Seguindo o mesmo padrão já usado para o \texttt{passwd} do Mosquitto: nenhuma credencial real em arquivo versionado. \texttt{infra/.env} (gitignored) guarda usuário, senha, organização, bucket e token; \texttt{infra/.env.example} é o template versionado. Uma \textit{fixture} \texttt{autouse} do pytest (\texttt{ensure\_env\_file}) garante que \texttt{.env} existe antes de qualquer teste, gerando valores aleatórios se ele não existir -- e nunca sobrescrevendo um \texttt{.env} real que o desenvolvedor já tenha configurado.

\subsection{Schema}

\begin{itemize}
    \item \textbf{Organização}: \texttt{fsae}. \textbf{Bucket}: \texttt{telemetria}.
    \item \textbf{Retenção}: infinita (sem limite configurado) -- decisão consciente, já que o objetivo é comparar treinos e estudar tendências ao longo de uma temporada inteira, não só o evento mais recente.
    \item \textbf{Measurement}: \texttt{telemetry}, já definido no checkpoint 4 (\texttt{backend/ingestion/transform.py}) -- um ponto por mensagem MQTT, os 20 sinais como campos.
\end{itemize}

\section{Testes}

Quatro testes cobrem este checkpoint especificamente (mais dois de checkpoints anteriores, estendidos):

\begin{enumerate}[label=\textbf{Teste \arabic*}]
    \item \texttt{test\_point\_written\_via\_production\_writer\_is\_queryable} -- escreve um ponto usando o \texttt{make\_influx\_writer} de produção (checkpoint 4) contra um InfluxDB efêmero real, e confirma via consulta Flux que o valor volta certo.
    \item \texttt{test\_multiple\_fields\_of\_same\_point\_are\_all\_queryable} -- confirma que múltiplos campos do mesmo ponto (\texttt{engineTemp}, \texttt{airTemp}) são gravados e consultados corretamente juntos.
    \item \texttt{test\_docker\_compose\_up\_starts\_a\_working\_influxdb} -- o \texttt{docker-compose.yml} real (não uma cópia da config) sobe um InfluxDB que responde no endpoint de saúde.
    \item \texttt{test\_full\_pipeline\_mock\_publisher\_to\_real\_influxdb} -- fecha o ciclo do caminho histórico inteiro: mock publisher $\to$ broker real $\to$ ingestão $\to$ InfluxDB real $\to$ consulta de volta, nenhum elo simulado.
\end{enumerate}

\subsection{Resultado da execução}

\begin{lstlisting}[caption={Saída real de pytest -v (arquivos deste checkpoint)}]
test_influxdb.py::test_point_written_via_production_writer_is_queryable PASSED
test_influxdb.py::test_multiple_fields_of_same_point_are_all_queryable PASSED
test_docker_compose.py::test_docker_compose_up_starts_a_working_influxdb PASSED
test_end_to_end.py::test_full_pipeline_mock_publisher_to_real_influxdb PASSED
\end{lstlisting}

\section{Descoberta durante a validação manual}

Depois de gerar um novo \texttt{infra/.env} (com um token de admin novo) e rodar \texttt{docker compose up}, uma tentativa de escrita real falhou com \texttt{401 Unauthorized}. Causa: o InfluxDB já tinha sido inicializado antes, numa sessão de teste anterior, e seu volume de dados (\texttt{influxdb-data}) persistiu entre execuções -- a imagem oficial só roda o processo de setup (que lê as variáveis de ambiente) na \textbf{primeira} subida com o volume vazio. Trocar \texttt{.env} depois não tem efeito nenhum até o volume ser apagado.

A correção foi \texttt{docker compose down -v} (remove os volumes) seguido de \texttt{up} novamente, forçando uma reinicialização completa. Esse comportamento foi documentado no \texttt{SCOPE.md} e no \texttt{CLAUDE.md} como nota operacional, pra não ser redescoberto por acidente mais pra frente.

\section{Nota sobre metodologia}

Diferente dos checkpoints anteriores, neste a implementação (serviço no \texttt{docker-compose.yml}, credenciais) foi escrita antes dos testes automatizados -- os testes vieram depois de uma validação manual já confirmar que tudo funcionava. Essa inversão de ordem foi identificada e discutida com o usuário após o checkpoint 6; a partir do checkpoint 7 o processo volta a ser teste-primeiro, inclusive para infraestrutura declarativa. A cobertura de testes deste checkpoint continua válida (roda contra infraestrutura real, não é simulada) -- só a ordem em que foi escrita que não seguiu TDD estrito.

\section{Resumo do que ficou pronto}

\begin{itemize}
    \item InfluxDB 2.x real, container Docker, setup automático via variáveis de ambiente.
    \item Credenciais nunca em arquivo versionado -- mesmo padrão do Mosquitto.
    \item Schema definido e documentado: organização, bucket, retenção infinita, measurement.
    \item Suposição de design do checkpoint 4 (escrita/leitura via \texttt{influxdb-client}) validada contra banco real, não mais mockada.
    \item Descoberta operacional sobre o comportamento de inicialização do InfluxDB documentada.
\end{itemize}

\section{Próximos passos}

Checkpoint 6: dashboards Grafana usando este InfluxDB como fonte de dados para o caminho histórico da arquitetura.

\end{document}
